% ===================================================
% CHAPITRE 14 — ALLOIGNET
% Pages imprimées 72-74 — vues 75-77
% ===================================================
\chapter{Alloignet}

% --- Page imprimée 72 — vue 75
\begin{center}\textit{[Illustration : ruines d'Alloignet]}\end{center}

La seigneurie d'Alloignet, située sur la paroisse de St-Mamert a joué, pendant plusieurs siècles un rôle trop important dans notre vallée pour que nous ne donnions pas une analyse succincte de son histoire.

Nous trouvons le nom d'Alloignet dès les premières années du XIIIe siècle. En Décembre 1224 Mathilde de Courtenay, Comtesse de Nivernais, tenait quitte Humbert de Beaujeu pour le fief de Couzan (Loire) en considération de ce qu'il a pris d'elle, en fief, Chevagny et Alloignet qu'il tenait précédemment en franc alleu. Le Seigneur de Beaujeu était en effet, pour le fief de Couzan, le vassal du Comte de Nevers, depuis qu'il avait cédé l'hommage de ce fief qu'il avait reçu du Comte Hugues de Damas.

Le nom d'Alloignet fut porté par une ancienne famille dont deux membres, Joannès et Stephanus de Alogneto furent religieux dans le célèbre monastère de l'Isle-Barbe et figurèrent au Chapitre en 1261

En 1281 Agnès, femme de Aymond Paladin, du consentement de son mari, vendit Alloignet à Louis Seigneur de Beaujeu. Le 15 des Calendes de Juin de l'an 1302 Guichard V donne à ses frères qui n'avaient reçu de leur père que 300 livres viennoises, plusieurs châteaux et entr'autres celui d'Alloignet.

Nous n'en retrouvons plus aucune trace jusqu'au commencement du XVIe siècle, époque où Philibert de Beaujeu, baron de Lignières et d'Amplepuis, ayant intenté un procès au duc de Bourbon en revendication des terres de Beaujolais et de Dombes, selon la substitution
% --- Page imprimée 73 — vue 76 (suite)
de Guichard V, reçut en vertu d'un arrangement en date du 5 Octobre 1516, les Seigneuries d'Alloignet, de Lay et d'Ussel, avec une rente de 1500 livres dont il se contenta pour tous droits.

La Châtellenie d'Alloignet comprenait les paroisses de St-Mamert d'Ouroux, St-Jacques, Trades, St-Christophe, les Ardillats et Chenelette.

Philibert de Beaujeu étant mort, sa veuve Catherine d'Amboise réclama la majeure partie de sa riche succession, par suite de la donation que son mari lui avait faite le 4 Février 1540, des terres d'Alloignet, etc... D'un autre côté la maison de Choiseul et celle de Barton de Montbar issues l'une et l'autre de Louis de Suilly époux de Marie de Beaujeu, sœur de Philibert, vinrent également en revendication des biens du baron de Lignières. Un premier arrêt du Parlement de Paris ordonna le partage de la succession. Mais l'affaire ayant été évoquée au Conseil du Roi et renvoyée au Parlement de Grenoble, la totalité des biens fut adjugée, par arrêt du 4 Juin 1573 à Charles Ludovic duc de Nevers qui, par contrat du 10 Mars 1578 en vendit la majeure partie à Claude de Rébé, seigneur du fief du même nom situé à Amplepuis.

Peu après, nous ne savons par suite de quelles circonstances, la seigneurie fit retour à la baronnie du Beaujolais possédée à cette époque par les Bourbon de Montpensier. Les titres qui concernent Alloignet et que nous avons reproduits au sujet des servis dus par les habitants d'Ouroux, citent fréquemment le nom de Mademoiselle, duchesse de Montpensier ou Anne Marie Louise d'Orléans. Ne résidant pas toujours dans le Beaujolais, ses biens furent toujours administrés par des tiers. Elle mourut en 1693 après avoir institué pour son héritier universel Philippe d'Orléans, frère unique du roi.

La châtellenie d'Alloignet fut toujours associée, dans ses diverses transformations, à la Prévôté de Coux. La seigneurie de Coux était située sur la paroisse de Monsols, au lieu appelé "le Chastelard".

En 1604 la place forte d'Alloignet n'était déjà plus qu'une vieille masure appelée le "Chasteau d'Alloignet" avec terre et bois taillis y attenant.

% --- Page imprimée 74 — vue 77
\begin{center}\textit{[Illustration : PLAN du bourg de Saint-Antoine d'Ouroux d'après l'ATLAS cadastral parcellaire terminé le 28 Octobre 1825 sous l'Administration de Mr le Comte de Brosses, préfet et de Mr Delaroche La Carelle fils, maire et sous la Direction de Mr Walle directeur des Contributions et de Mr Gaud, géomètre en chef, par Mr Terra géomètre du cadastre.]}\end{center}
